\section{Why this matters}\label{sec:intro}

Language models are entering retail financial advice and research workflows, and central banks have begun to map how AI changes intermediation and asset management~\cite{aldasoro2024,bis2024ai}. If a model systematically allocated more to one country's assets, that tilt would reach every user who asks it, much as investor home bias~\cite{french1991,lewis1999,coval1999} shapes portfolios one investor at a time. Language models also carry geographic and cultural slants: their default answers resemble the values of English-speaking and Protestant European countries~\cite{tao2024,durmus2023}, and their ratings of places are systematically lower for regions with weaker socioeconomic conditions~\cite{manvi2024}. In finance, a recent audit of 32 models finds advice consistent with investor home bias~\cite{fieberg2025}, and three commercial chat assistants put over 90\% of equity portfolios in the United States~\cite{winder2025}; whether advice is trustworthy enough to be fiduciary is open~\cite{lo2024}.

Those audits vary the investor, not the model's origin. This note asks the complementary question with open-weight models whose developers sit in different countries: \emph{holding the prompt fixed, does a model allocate more to its own home market?} Because the prompt also states the investor's residence and is written in four languages, the design separates three channels that are confounded in ordinary use: the model's origin~($\beta$), the investor's residence~($\gamma$) and the prompt language~($\lambda$). Two Thai models fine-tuned from American and Chinese base models that are also in the sample let us ask whether fine-tuning moves the allocation towards Thailand.
